\section{Exercises}
\label{sec:A.8}

\begin{enumerate}

\item[A.1] You run a Markov chain on the set $\{1, 2\}$ with transition matrix
\[
  P = \frac{1}{3}\begin{bmatrix} 2 & 1 \\ 1 & 2 \end{bmatrix}.
\]
For general positive integer $n$, calculate $P^n$.  Give your answer as simply as possible.

\item[A.2] Let $E = \{1, \ldots, s\}$ with $s$ a positive integer.  We consider Markov chains
starting from $X_0 = 1$.  Find transition matrices $P : E \times E \to \R$ that generate Markov
chains with the following properties.  If no such $P$ exists, explain why not.

\begin{enumerate}[label=\alph*)]
  \item The chain is periodic with cycle~$3$.

  \item For all $n$, $X_{n+1} = X_n$ with probability $1/2$, and the chain is periodic with
    cycle~$4$.

  \item The size of the state space is $s = 2$ and the chain reaches equilibrium in one step.

  \item The chain is irreducible and aperiodic, but after $10^{12}$ steps the chain will not be
    close to its equilibrium distribution.
\end{enumerate}

\item[A.3] This exercise illustrates that the detailed balance equations are often easier to solve
than the general balance equations.  Consider the weather chain example with transition matrix
\[
  P = \begin{bmatrix} 0.4 & 0.6 & 0 \\ 0.25 & 0.25 & 0.5 \\ 0 & 0.4 & 0.6 \end{bmatrix}.
\]
Write and solve the general balance equation.  Write and solve the detailed balance equations.

\item[A.4] Consider a Markov chain on the set $\{-1, 1\}$ with transition matrix
\[
  P = \begin{bmatrix} \alpha & 1 - \alpha \\ 1 - \alpha & \alpha \end{bmatrix}.
\]
Find a simple formula for its autocorrelation of lag~$k$.  As a sanity check, make sure that
your answer makes sense for the specific values $\alpha = 0$, $\alpha = 1/2$, and $\alpha = 1$.

\end{enumerate}
